\section{Transición areal--volumétrica y escala espectral}
\label{sec:ym-areal-volumetric-scale}

La constante que fija el semigrupo gauge se obtiene antes de introducir el
Hamiltoniano. Su procedencia combina dos lectores independientes del mismo
estado genealógico: la frecuencia de retorno que expresa
\(\pi_{\rm HMT}\) y la distribución estacionaria de las dos hojas
areal--volumétrica.

\subsection{Dinámica de las dos hojas}

Sea
\begin{equation}
 F_{\rm av}=\begin{pmatrix}0&1\\1&1\end{pmatrix},
 \qquad e_{\rm ar}=\binom10,
 \qquad e_{\rm vol}=\binom01.
 \label{eq:ym-av-transition-matrix}
\end{equation}
La recurrencia de Fibonacci da, por inducción,
\begin{equation}
 F_{\rm av}^{n}
 =\begin{pmatrix}F_{n-1}&F_n\\F_n&F_{n+1}\end{pmatrix}.
 \label{eq:ym-av-fibonacci-powers}
\end{equation}
El vector de Perron es proporcional a \((1,\varphi)\). Como
\(F_{\rm av}\) es simétrica, la medida de Parry asigna a las dos hojas
pesos proporcionales a \((1,\varphi^2)\), y por tanto
\begin{equation}
 \frac{\mu_{\rm ar}}{\mu_{\rm vol}}=\varphi^{-2},
 \qquad
 \mu_{\rm vol}=\frac{\varphi^2}{1+\varphi^2}.
 \label{eq:ym-av-parry-weights}
\end{equation}
Aquí \(\mu_{\rm vol}\) es un carácter de frecuencia adimensional; no
denota un volumen de dimensión \(L^4\).

El lector regional de \(\pi\) se introduce una sola vez mediante
\begin{equation}
 D_\pi
 =\pi_{\rm HMT}|e_{\rm ar}\rangle\langle e_{\rm ar}|
  +|e_{\rm vol}\rangle\langle e_{\rm vol}|.
 \label{eq:ym-pi-regional-reader}
\end{equation}
La razón expresada a profundidad \(n\) es
\begin{equation}
 \Theta_n
 :=\frac{\langle e_{\rm ar},D_\pi F_{\rm av}^{n}e_{\rm ar}\rangle}
 {\langle e_{\rm vol},F_{\rm av}^{n}e_{\rm vol}\rangle}
 =\pi_{\rm HMT}\frac{F_{n-1}}{F_{n+1}}
 \longrightarrow
 r_{\rm av}:=\frac{\pi_{\rm HMT}}{\varphi_{\rm HMT}^{2}}.
 \label{eq:ym-target-free-av-limit}
\end{equation}
La matriz, la ruta y el cociente se fijan antes de consultar un salto de
masa o un parámetro de teoría gauge. En la profundidad nativa \(n=108\),
\begin{equation}
 F_{107}=10284720757613717413913,
 \qquad
 F_{109}=26925748508234281076009,
 \label{eq:ym-fibonacci-107-109}
\end{equation}
y
\begin{equation}
 \left|\Theta_{108}
 -\frac{\pi_{\rm HMT}}{\varphi_{\rm HMT}^{2}}\right|
 <2\cdot10^{-45}.
 \label{eq:ym-av-depth-108-error}
\end{equation}

\subsection{Transducción dimensional}

Sea \(\ell_0\) la unidad longitudinal del lector dimensional. La escala
inversa de longitud y el parámetro de una losa de espesor \(a_m\) son
\begin{equation}
 \kappa_{\rm av}
 :=\frac{\mu_{\rm vol}}{\ell_0}
 \left(\frac{\pi_{\rm HMT}}{\varphi_{\rm HMT}^{2}}-1\right)>0,
 \qquad
 q_m:=e^{-3a_m\kappa_{\rm av}}.
 \label{eq:ym-dimensional-gap-scale}
\end{equation}
Por tanto \([\kappa_{\rm av}]=L^{-1}\) y la brecha energética es
\begin{equation}
 \Delta_E=3\hbar c\,\kappa_{\rm av};
 \label{eq:ym-energy-gap-units}
\end{equation}
en unidades \(\hbar=c=1\), se escribe \(\Delta_E=3\kappa_{\rm av}\).
La memoria electrónica \(\Delta_4\) pertenece a otro dominio y a otro
lector; no interviene en \eqref{eq:ym-dimensional-gap-scale}.

En los semigrupos de las secciones siguientes, \(t,s,a_m\) son
longitudes euclídeas y los generadores \(D\) y \(H^{\rm ov}\) tienen
dimensión \(L^{-1}\). El Hamiltoniano energético es
\(H_E=\hbar_{\rm int}c_{\rm int}D\). Si
\(\tau=t/c_{\rm int}\), entonces
\[
 e^{-tD}=e^{-\tau H_E/\hbar_{\rm int}},\qquad
 \Delta=3\kappa_{\rm av},\quad
 \Delta_E=\hbar_{\rm int}c_{\rm int}\Delta,\quad
 m_{\rm gap}=\frac{\hbar_{\rm int}\Delta}{c_{\rm int}}.
\]
La misma brecha queda así expresada, respectivamente, como longitud
inversa, energía y masa en la coordenatización dimensional declarada.

Como \(a_m=9a_{m+1}\), la definición da inmediatamente
\begin{equation}
 q_{m+1}^{9}=q_m.
 \label{eq:ym-nonadic-gap-root}
\end{equation}
Esta identidad es la raíz nonádica que transportará el mismo semigrupo y la
misma escala espectral a través de todos los refinamientos.
